% Lösungen Punkte und Strecken im Koordinatensystem (zusammenbau v0.4, Heft)
\einheitenkopf{Punkte und Strecken im Koordinatensystem}

\erg{1}{a) $P$ liegt in der $x_1x_2$-Ebene \quad b) $C(4 | 5 | 2)$; $A$ und $B$ siehe Lösungsgrafik \quad c) Spitzen $(-2 | -2 | 4)$ und $(5 | 3 | 3)$; der erste Mast steht hinter der $x_2x_3$-Ebene und links der $x_1x_3$-Ebene; siehe Lösungsgrafik.}

1b)

\begin{ksys3} \rpunkt{5}{3}{1}{A} \rpunkt{2}{1}{4}{B} \rpunkt{4}{5}{2}{C} \end{ksys3}


1c)

\begin{ksys3} \rpunkt*{-2}{-2}{0}{F_1} \rpunkt*{5}{3}{0}{F_2} \rgerade[0:1]{-2,-2,0}{0,0,4}{} \rgerade[0:1]{5,3,0}{0,0,3}{} \rgerade[0:1]{-2,-2,4}{7,5,-1}{} \end{ksys3}

\erg{2}{a) Jeder Punkt hat zwei Koordinaten $0$ oder $4$ und liegt damit auf einer Würfelkante; die vier Punkte der Reihe nach verbinden, siehe Lösungsgrafik. \quad b) Jeder Punkt hat zwei Koordinaten $0$ oder $3$ und liegt damit auf einer Würfelkante; die vier Punkte der Reihe nach verbinden, siehe Lösungsgrafik. \quad c) Jeder Punkt hat zwei Koordinaten $0$ oder $4$ und liegt damit auf einer Würfelkante; die vier Punkte der Reihe nach verbinden, siehe Lösungsgrafik. \quad d) $P$ liegt auf der Deckkante mit $x_1 = 5$, ein Drittel ihrer Länge von vorn; Dreieck siehe Lösungsgrafik. \quad e) $P$ liegt auf der senkrechten Kante über $(6 | 0 | 0)$, ein Viertel der Höhe; Dreieck siehe Lösungsgrafik. \quad f) $P$ liegt auf der Deckkante mit $x_1 = 3$, ein Drittel ihrer Länge von vorn; Dreieck siehe Lösungsgrafik.}

2a)

\begin{ksys3} \rquader{0,0,0}{4}{4}{4} \rpunkt*{4}{0}{1}{P} \rpunkt*{1}{4}{0}{Q} \rpunkt*{0}{4}{2}{R} \rpunkt*{3}{0}{4}{S} \rgerade[0:1]{4,0,1}{-3,4,-1}{} \rgerade[0:1]{1,4,0}{-1,0,2}{} \rgerade[0:1]{0,4,2}{3,-4,2}{} \rgerade[0:1]{3,0,4}{1,0,-3}{} \end{ksys3}


2b)

\begin{ksys3} \rquader{0,0,0}{3}{3}{3} \rpunkt*{3}{1}{0}{P} \rpunkt*{3}{3}{2}{Q} \rpunkt*{1}{3}{3}{R} \rpunkt*{0}{0}{2}{S} \rgerade[0:1]{3,1,0}{0,2,2}{} \rgerade[0:1]{3,3,2}{-2,0,1}{} \rgerade[0:1]{1,3,3}{-1,-3,-1}{} \rgerade[0:1]{0,0,2}{3,1,-2}{} \end{ksys3}


2c)

\begin{ksys3} \rquader{0,0,0}{4}{4}{4} \rpunkt*{0}{3}{0}{P} \rpunkt*{4}{4}{1}{Q} \rpunkt*{4}{1}{4}{R} \rpunkt*{1}{0}{4}{S} \rgerade[0:1]{0,3,0}{4,1,1}{} \rgerade[0:1]{4,4,1}{0,-3,3}{} \rgerade[0:1]{4,1,4}{-3,-1,0}{} \rgerade[0:1]{1,0,4}{-1,3,-4}{} \end{ksys3}


2d)

\begin{ksys3} \rquader{0,0,0}{5}{6}{3} \rpunkt*{0}{0}{0}{O} \rpunkt*{5}{2}{3}{P} \rpunkt*{0}{6}{3}{R} \rgerade[0:1]{0,0,0}{5,2,3}{} \rgerade[0:1]{5,2,3}{-5,4,0}{} \rgerade[0:1]{0,6,3}{0,-6,-3}{} \end{ksys3}


2e)

\begin{ksys3} \rquader{0,0,0}{6}{3}{4} \rpunkt*{0}{0}{0}{O} \rpunkt*{6}{0}{1}{P} \rpunkt*{0}{3}{4}{T} \rgerade[0:1]{0,0,0}{6,0,1}{} \rgerade[0:1]{6,0,1}{-6,3,3}{} \rgerade[0:1]{0,3,4}{0,-3,-4}{} \end{ksys3}


2f)

\begin{ksys3} \rquader{0,0,0}{3}{6}{3} \rpunkt*{0}{0}{0}{O} \rpunkt*{3}{2}{3}{P} \rpunkt*{3}{6}{0}{T} \rgerade[0:1]{0,0,0}{3,2,3}{} \rgerade[0:1]{3,2,3}{0,4,-3}{} \rgerade[0:1]{3,6,0}{-3,-6,0}{} \end{ksys3}

\erg{3}{a) Die dritten Koordinaten $3$ und $2$ sind beide positiv: beide Punkte liegen oberhalb der $x_1x_2$-Ebene. \quad b) Alle Ecken haben $x_1 = 6$: parallel zur $x_2x_3$-Ebene. $B$ und $C$ unterscheiden sich nur im Vorzeichen von $x_2$, $A$ hat $x_2 = 0$: symmetrisch zur $x_1x_3$-Ebene. \quad c) Die dritten Koordinaten $-3$ und $2$ haben verschiedene Vorzeichen: $R$ liegt unterhalb, $T$ oberhalb der $x_1x_2$-Ebene. \quad d) Fehlende Ecken $(-5 | 3)$, $(-6 | 0)$, $(-5 | -3)$, $(0 | -4)$, $(5 | -3)$; siehe Lösungsgrafik. \quad e) Fehlende Ecken $(-3 | 2)$, $(-4 | 0)$, $(-3 | -2)$, $(0 | -3)$, $(3 | -2)$; siehe Lösungsgrafik. \quad f) Fehlende Ecken $(-2 | 5)$, $(-6 | 0)$, $(-2 | -5)$, $(0 | -6)$, $(2 | -5)$; siehe Lösungsgrafik.}

3d)

\begin{ksys}[xmin=-7,xmax=7,ymin=-7,ymax=7,klein] \punkt{6}{0}{} \punkt{5}{3}{} \punkt{0}{4}{} \punkt{-5}{3}{} \punkt{-6}{0}{} \punkt{-5}{-3}{} \punkt{0}{-4}{} \punkt{5}{-3}{} \end{ksys}


3e)

\begin{ksys}[xmin=-7,xmax=7,ymin=-7,ymax=7,klein] \punkt{4}{0}{} \punkt{3}{2}{} \punkt{0}{3}{} \punkt{-3}{2}{} \punkt{-4}{0}{} \punkt{-3}{-2}{} \punkt{0}{-3}{} \punkt{3}{-2}{} \end{ksys}


3f)

\begin{ksys}[xmin=-7,xmax=7,ymin=-7,ymax=7,klein] \punkt{6}{0}{} \punkt{2}{5}{} \punkt{0}{6}{} \punkt{-2}{5}{} \punkt{-6}{0}{} \punkt{-2}{-5}{} \punkt{0}{-6}{} \punkt{2}{-5}{} \end{ksys}

\erg{4}{a) $|SB| = |SD| = \sqrt{9 + 16} = 5$ cm (nicht $4$ cm); Spitzen bei $(-4 | 0)$ an $AD$, $(8 | 0)$ an $BC$ und $(0 | 8)$ an $CD$ – die Dreiecke an $BC$ und $CD$ sind bei $B$ bzw. $D$ rechtwinklig. \quad b) $|SB| = |SD| = \sqrt{16 + 9} = 5$ cm (nicht $3$ cm); Spitzen bei $(-3 | 0)$ an $AD$, $(9 | 0)$ an $BC$ und $(0 | 9)$ an $CD$ – die Dreiecke an $BC$ und $CD$ sind bei $B$ bzw. $D$ rechtwinklig. \quad c) $|SB| = |SD| = \sqrt{4 + 2{,}25} = 2{,}5$ cm (nicht $1{,}5$ cm); Spitzen bei $(-1{,}5 | 0)$ an $AD$, $(4{,}5 | 0)$ an $BC$ und $(0 | 4{,}5)$ an $CD$ – die Dreiecke an $BC$ und $CD$ sind bei $B$ bzw. $D$ rechtwinklig. \quad d) Die Grundfläche liegt parallel zur $x_1x_2$-Ebene, also zeigt $S'$ die Lage des Höhenfußpunkts: $C'(-1 | -1 | 0)$; $S'(4 | -1 | 0)$ liegt außerhalb des Dreiecks $A'B'C'$ – der Höhenfußpunkt liegt außerhalb der Grundfläche. \quad e) Die Grundfläche liegt parallel zur $x_1x_2$-Ebene, also zeigt $S'$ die Lage des Höhenfußpunkts: $C'(-1 | -2 | 0)$; $S'(1 | 1 | 0)$ liegt innerhalb des Dreiecks $A'B'C'$ – der Höhenfußpunkt liegt innerhalb der Grundfläche.}

4a)

\begin{ksys}[xmin=-5,xmax=9,ymin=-5,ymax=9,klein] \punkt{0}{0}{A} \punkt{3}{0}{B} \punkt{3}{3}{C} \punkt{0}{3}{D} \punkt{0}{-4}{S} \punkt{-4}{0}{S} \punkt{8}{0}{S} \punkt{0}{8}{S} \end{ksys}


4b)

\begin{ksys}[xmin=-4,xmax=10,ymin=-4,ymax=10,klein] \punkt{0}{0}{A} \punkt{4}{0}{B} \punkt{4}{4}{C} \punkt{0}{4}{D} \punkt{0}{-3}{S} \punkt{-3}{0}{S} \punkt{9}{0}{S} \punkt{0}{9}{S} \end{ksys}


4c)

\begin{ksys}[xmin=-2.5,xmax=5.5,ymin=-2.5,ymax=5.5,klein] \punkt{0}{0}{A} \punkt{2}{0}{B} \punkt{2}{2}{C} \punkt{0}{2}{D} \punkt{0}{-1.5}{S} \punkt{-1.5}{0}{S} \punkt{4.5}{0}{S} \punkt{0}{4.5}{S} \end{ksys}


4d)

\begin{ksys3} \rpunkt*{4}{1}{0}{A'} \rpunkt*{1}{5}{0}{B'} \rpunkt*{-1}{-1}{0}{C'} \rpunkt*{4}{-1}{0}{S'} \rgerade[0:1]{4,1,0}{-3,4,0}{} \rgerade[0:1]{1,5,0}{-2,-6,0}{} \rgerade[0:1]{-1,-1,0}{5,2,0}{} \end{ksys3}


4e)

\begin{ksys3} \rpunkt*{5}{-1}{0}{A'} \rpunkt*{1}{5}{0}{B'} \rpunkt*{-1}{-2}{0}{C'} \rpunkt*{1}{1}{0}{S'} \rgerade[0:1]{5,-1,0}{-4,6,0}{} \rgerade[0:1]{1,5,0}{-2,-7,0}{} \rgerade[0:1]{-1,-2,0}{6,1,0}{} \end{ksys3}

\erg{5}{a) ein Punkt – die Mitte von $AB$ \quad b) $\overrightarrow{AB} = (2 | 3 | 6)$, $|\overrightarrow{AB}| = \sqrt{4 + 9 + 36} = 7$ LE \quad c) $|\overrightarrow{SH}| = \sqrt{4 + 9 + 36} = 7$ m; mit Zuschlag $0{,}4$ m (nicht $40$ m): $4 \cdot 7{,}4 = 29{,}6$ m}
\erg{6}{a) $|\overrightarrow{AB}| = \sqrt{36 + 64 + 0} = 10$ m; vorher $10 : 1{,}25 = 8$ m (nicht $10$ minus $25\,\%$) \quad b) $|\overrightarrow{AB}| = \sqrt{4 + 36 + 81} = 11$ m; vorher $11 : 1{,}1 = 10$ m (nicht $11$ minus $10\,\%$) \quad c) $|\overrightarrow{AB}| = \sqrt{1 + 16 + 64} = 9$ cm; vorher $9 : 1{,}5 = 6$ cm (nicht $9$ minus $50\,\%$) \quad d) $\frac12 \cdot (\overrightarrow{OD} + \overrightarrow{OS}) = (2 | 1{,}5 | 3)$, also ist $W$ Mittelpunkt; an der $x_1x_3$-Ebene gespiegelt: $T(2 | -1{,}5 | 3)$ \quad e) $\frac12 \cdot (\overrightarrow{OD} + \overrightarrow{OS}) = (3 | 3 | 4)$, also ist $W$ Mittelpunkt; an der $x_2x_3$-Ebene gespiegelt: $T(-3 | 3 | 4)$ \quad f) $\frac12 \cdot (\overrightarrow{OD} + \overrightarrow{OS}) = (4 | 3 | 2)$, also ist $W$ Mittelpunkt; an der $x_1x_2$-Ebene gespiegelt: $T(4 | 3 | -2)$ \quad g) $T_1 = P + \frac{1}{3} \cdot \overrightarrow{PQ} = (6 | 1 | 2)$, $T_2 = P + \frac{2}{3} \cdot \overrightarrow{PQ} = (3 | 2 | 1)$ (nicht die Mitte) \quad h) Anteil $\frac{2}{5}$ (nicht $\frac{2}{3}$): $G = A + \frac{2}{5} \cdot \overrightarrow{AC} = (5 | 2 | 4)$ \quad i) Anteil $\frac{3}{4}$ (nicht $\frac{1}{4}$ – $G$ liegt näher an $C$): $G = A + \frac{3}{4} \cdot \overrightarrow{AC} = (5 | 6 | 2)$ \quad j) $D = A - \overrightarrow{AB} = (1 | 4 | 0)$ (nicht $B$ und nicht $A + 2\overrightarrow{AB}$) \quad k) $D = A - \overrightarrow{AB} = (-2 | -6 | 4)$ (nicht $B$ und nicht $A + 2\overrightarrow{AB}$) \quad l) $D = A - \overrightarrow{AB} = (-2 | 5 | 4)$ (nicht $B$ und nicht $A + 2\overrightarrow{AB}$)}
\erg{7}{a) $\overrightarrow{AC} = \overrightarrow{OC} - \overrightarrow{OA} = 2 \cdot \overrightarrow{OA} = (4 | -2 | 4)$, $|\overrightarrow{AC}| = 6$ (nicht $|\overrightarrow{OC}|$) \quad b) $\overrightarrow{AC} = \overrightarrow{OC} - \overrightarrow{OA} = -2 \cdot \overrightarrow{OA} = (-4 | -6 | -12)$, $|\overrightarrow{AC}| = 14$ (nicht $|\overrightarrow{OC}|$) \quad c) $\overrightarrow{AC} = \overrightarrow{OC} - \overrightarrow{OA} = -\frac{1}{2} \cdot \overrightarrow{OA} = (-2 | -2 | -1)$, $|\overrightarrow{AC}| = 3$ (nicht $|\overrightarrow{OC}|$) \quad d) Tunnellänge $|\overrightarrow{AB}| \cdot 10$ m $= 500$ m; gleiche Fahrzeit: $\frac{s}{90} = \frac{500 - s}{60}$ ergibt $s = 300$ m (Weg des Zugs ab $B$). \quad e) $|\overrightarrow{PQ}| = 10$, also $1000$ m; $\frac{s}{15} = \frac{1000 - s}{10}$ ergibt $s = 600$ m; Treffpunkt $P + 0{,}6 \cdot \overrightarrow{PQ} = (5{,}6 | 5{,}8 | 0)$}
\erg{8}{a) $\overrightarrow{DC}$}
\erg{9}{a) $|\overrightarrow{CA}| = |\overrightarrow{CB}| = \sqrt{21}$, $|\overrightarrow{AB}| = \sqrt{20}$: gleichschenklig, Schenkel $CA$ und $CB$, Basis $AB$ \quad b) $|\overrightarrow{CA}| = |\overrightarrow{CB}| = \sqrt{14}$, $|\overrightarrow{AB}| = \sqrt{20}$: gleichschenklig, Schenkel $CA$ und $CB$, Basis $AB$ \quad c) $|\overrightarrow{RP}| = |\overrightarrow{RQ}| = \sqrt{29}$, $|\overrightarrow{PQ}| = \sqrt{20}$: gleichschenklig, Schenkel $RP$ und $RQ$, Basis $PQ$ \quad d) $|\overrightarrow{CA}| = |\overrightarrow{CB}| = \sqrt{14}$, $|\overrightarrow{AB}| = \sqrt{24}$: gleichschenklig, Schenkel $CA$ und $CB$, Basis $AB$ \quad e) $|\overrightarrow{CA}| = |\overrightarrow{CB}| = \sqrt{29}$, $|\overrightarrow{AB}| = 6$: gleichschenklig, Schenkel $CA$ und $CB$, Basis $AB$ \quad f) $|\overrightarrow{AB}| = |\overrightarrow{BC}| = |\overrightarrow{CA}| = \sqrt{8}$: gleichseitig. \quad g) $|\overrightarrow{AB}| = \sqrt{18}$, $|\overrightarrow{BC}| = |\overrightarrow{CA}| = 5$: nicht gleichseitig, nur gleichschenklig mit der Basis $AB$. \quad h) $|\overrightarrow{AB}| = |\overrightarrow{BC}| = |\overrightarrow{CA}| = \sqrt{32}$: gleichseitig. \quad i) $|\overrightarrow{CA}| = |(-2 | -2 | -p)| = \sqrt{8 + p^2}$ und $|\overrightarrow{CB}| = |(2 | 2 | -p)| = \sqrt{8 + p^2}$: gleich für jedes $p$, Basis $AB$. \quad j) $|\overrightarrow{CA}| = |\overrightarrow{CB}| = \sqrt{(3 - p)^2 + 10}$, denn $\overrightarrow{CA} = (3 - p | -3 | 1)$ und $\overrightarrow{CB} = (3 - p | 3 | -1)$ unterscheiden sich nur im Vorzeichen zweier Koordinaten. \quad k) $|\overrightarrow{CA}| = |\overrightarrow{CB}| = \sqrt{8 + (2 - p)^2}$ für jedes $p$; Basis $AB$. \quad l) Mit dem Ursprung $O$ ist $OABC$ ein Quadrat (Seiten auf den Achsen bzw. parallel dazu, alle gleich lang); $ABC$ ist seine Hälfte mit dem rechten Winkel bei $B$ (nicht bei $A$ oder $C$) und $|AB| = |BC|$. Flächeninhalt $\frac12 \cdot 9 \cdot 9 = 40{,}5$.}
\erg{10}{a) $|\overrightarrow{AB}| = |\overrightarrow{AC}| = \sqrt{34}$, Basis $BC$ (nicht Schenkel); $E$ und $F$ haben dieselbe dritte Koordinate $2$, die Grundfläche liegt in der $x_1x_2$-Ebene: $EF$ ist parallel zu ihr. \quad b) $|\overrightarrow{AB}| = |\overrightarrow{AC}| = 5$, Basis $BC$ (nicht Schenkel); $G$ und $H$ haben dieselbe dritte Koordinate $5$, die Grundfläche liegt in der Ebene $x_3 = 3$: $GH$ ist parallel zu ihr. \quad c) $|\overrightarrow{AB}| = |\overrightarrow{AC}| = \sqrt{29}$, Basis $BC$ (nicht Schenkel); $P$ und $Q$ haben dieselbe dritte Koordinate $3$, die Grundfläche liegt in der $x_1x_2$-Ebene: $PQ$ ist parallel zu ihr. \quad d) $M$ ist der Mittelpunkt von $DE$; $\overrightarrow{FD} \circ \overrightarrow{FE} = 0$, also hat das Dreieck bei $F$ einen rechten Winkel. Nach Thales liegt $F$ auf dem Kreis um $M$ durch $D$ und $E$. \quad e) $\overrightarrow{FD} \circ \overrightarrow{FE} = 0$: rechter Winkel bei $F$. Nach Thales liegt $F$ auf dem Kreis mit dem Durchmesser $DE$, sein Mittelpunkt ist die Mitte von $DE$ – drei Abstände muss man nicht rechnen. \quad f) $\overrightarrow{RP} \circ \overrightarrow{RQ} = 0$: rechter Winkel bei $R$. Nach Thales liegt $R$ auf dem Kreis mit dem Durchmesser $PQ$, sein Mittelpunkt ist die Mitte von $PQ$ – drei Abstände muss man nicht rechnen. \quad g) Symmetrie zu $AB$: $M(3 | m)$; $|MA| = |MC|$: $3^2 + m^2 = (9 - m)^2$ ergibt $m = 4$: $M(3 | 4)$ (nicht der Schwerpunkt $(3 | 3)$) \quad h) Symmetrie zu $AB$: $M(2 | m)$; $|MA| = |MC|$: $4^2 + m^2 = (8 - m)^2$ ergibt $m = 3$: $M(2 | 3)$ \quad i) Symmetrie zu $AB$: $M(4 | m)$; $|MA| = |MC|$: $3^2 + (m - 1)^2 = (10 - m)^2$ ergibt $m = 5$: $M(4 | 5)$}
\erg{11}{a) $C$ am Mittelpunkt $M(3 | 3 | 3)$ von $AB$ spiegeln: $D = A + B - C = (1 | 3 | 5)$; $ADBC$ ist dann eine Raute. \quad b) $C$ am Mittelpunkt $M(2 | 2 | 1)$ von $AB$ spiegeln: $D = A + B - C = (0 | 1 | -1)$; $ADBC$ ist dann eine Raute. \quad c) $|\overrightarrow{MA}| = 3$; $B = 2M - A = (-1 | 5 | 0)$ (Gegenpunkt); $C = M + (1 | 2 | 2) = (2 | 5 | 3)$, denn $(1 | 2 | 2)$ steht senkrecht auf $\overrightarrow{MA}$ und hat die Länge $3$ – dann ist $|CA| = |CB|$. \quad d) $|\overrightarrow{MA}| = 3$; $B = 2M - A = (6 | 4 | 3)$ (Gegenpunkt); $C = M + (1 | -2 | -2) = (5 | 0 | 2)$, denn $(1 | -2 | -2)$ steht senkrecht auf $\overrightarrow{MA}$ und hat die Länge $3$ – dann ist $|CA| = |CB|$.}
\erg{12}{a) Parallelogramm und ein rechter Winkel}
\erg{13}{a) $\overrightarrow{AB} = (2 | 2 | 1) = \overrightarrow{DC}$: Parallelogramm; $\overrightarrow{AB} \circ \overrightarrow{AD} = (2 | 2 | 1) \circ (-2 | 1 | 2) = 0$: rechter Winkel, also Rechteck. \quad b) $\overrightarrow{AB} = (3 | 0 | 0) = \overrightarrow{DC}$: Parallelogramm; $\overrightarrow{BC} = (0 | -4 | 5 - p)$, $\overrightarrow{AB} \circ \overrightarrow{BC} = 0$ für jedes $p$: Rechteck; $|\overrightarrow{BC}| = \sqrt{16 + (5 - p)^2}$. \quad c) $\overrightarrow{AB} = \overrightarrow{DC} = (2 | 2 | 1)$: Parallelogramm; $|\overrightarrow{AB}| = |\overrightarrow{AD}| = 3$: alle vier Seiten gleich lang (nicht nur zwei), also Raute. \quad d) $\overrightarrow{AB} = \overrightarrow{DC} = (4 | 0 | 3)$: Parallelogramm; $|\overrightarrow{AB}| = |\overrightarrow{AD}| = 5$: alle vier Seiten gleich lang (nicht nur zwei), also Raute. \quad e) $\overrightarrow{AB} = \overrightarrow{DC} = (1 | 4 | 8)$: Parallelogramm; $|\overrightarrow{AB}| = |\overrightarrow{AD}| = 9$: alle vier Seiten gleich lang (nicht nur zwei), also Raute. \quad f) $\overrightarrow{AB} = (3 | 1 | 0) = \overrightarrow{DC}$: Parallelogramm; $\overrightarrow{AB} \circ \overrightarrow{AD} = 5 \ne 0$: kein rechter Winkel bei $A$, also kein Rechteck. \quad g) $\overrightarrow{AB} = \overrightarrow{DC} = (2 | 1 | 2)$ und $|\overrightarrow{AB}| = |\overrightarrow{AD}| = 3$: Raute; $\overrightarrow{AB} \circ \overrightarrow{AD} = 4 \ne 0$: kein Quadrat. \quad h) $\overrightarrow{AB} = \overrightarrow{DC} = (0 | 3 | 4)$ und $|\overrightarrow{AB}| = |\overrightarrow{AD}| = 5$: Raute; $\overrightarrow{AB} \circ \overrightarrow{AD} = 12 \ne 0$: kein Quadrat. \quad i) $\overrightarrow{EF} = (-6 | 6 | 0) = 2 \cdot \overrightarrow{AB}$: $AB \parallel EF$, also Trapez (nicht aus gleichen Längen schließen); $|\overrightarrow{AE}| = |\overrightarrow{BF}| = \sqrt{18}$. \quad j) $\overrightarrow{PS} = (-6 | 0 | 6) = 2 \cdot \overrightarrow{QR}$: $PS \parallel QR$ und doppelt so lang; $|\overrightarrow{PQ}| = |\overrightarrow{SR}| = \sqrt{93}$.}
\erg{14}{a) $\overrightarrow{AD} = (0 | 6 | 0) = 3 \cdot \overrightarrow{BC}$: Trapez. Höhe ist der Abstand der Parallelen in $x_1$-Richtung, $h = 6$ (nicht $|AB|$). Flächeninhalt $\frac12 \cdot (6 + 2) \cdot 6 = 24$ m². \quad b) $\overrightarrow{EF} = (8 | 0 | 0) = 4 \cdot \overrightarrow{HG}$: Trapez. Höhe $h = 4$ (Abstand in $x_2$-Richtung, nicht der Schenkel $|FG| = 5$). Flächeninhalt $\frac12 \cdot (8 + 2) \cdot 4 = 20$ m². \quad c) $\overrightarrow{KN} = (0 | 6 | 0) = 3 \cdot \overrightarrow{LM}$: Trapez. Höhe $h = 4$ (Abstand in $x_1$-Richtung). Flächeninhalt $\frac12 \cdot (6 + 2) \cdot 4 = 16$ dm². \quad d) $|\overrightarrow{KL}| = |\overrightarrow{KN}| = \sqrt{13}$ und $|\overrightarrow{ML}| = |\overrightarrow{MN}| = 5$: zwei Paare benachbarter gleich langer Seiten (die Gegenseiten sind verschieden lang). \quad e) $|\overrightarrow{KL}| = |\overrightarrow{KN}| = \sqrt{13}$ und $|\overrightarrow{ML}| = |\overrightarrow{MN}| = \sqrt{45}$: zwei Paare benachbarter gleich langer Seiten (die Gegenseiten sind verschieden lang). \quad f) $|\overrightarrow{PQ}| = |\overrightarrow{PS}| = 4$ und $|\overrightarrow{RQ}| = |\overrightarrow{RS}| = \sqrt{10}$: zwei Paare benachbarter gleich langer Seiten (die Gegenseiten sind verschieden lang). \quad g) $\overrightarrow{AE} = (0 | 6 | 1) = 2 \cdot \overrightarrow{CD}$ mit $\overrightarrow{CD} = (0 | 3 | 0{,}5)$: parallel; $\overrightarrow{CD} \circ \overrightarrow{DE} = (0 | 3 | 0{,}5) \circ (-6 | 0 | 0) = 0$: rechter Winkel bei $D$. \quad h) $\overrightarrow{AB} = (2 | 2 | 1) = 2 \cdot \overrightarrow{DC}$ mit $\overrightarrow{DC} = (1 | 1 | 0{,}5)$: parallel; $\overrightarrow{AB} \circ \overrightarrow{AD} = (2 | 2 | 1) \circ (-1 | 2 | -2) = 0$: rechter Winkel bei $A$. \quad i) $\overrightarrow{AB} = (6 | 0 | 8) = 2 \cdot \overrightarrow{DC}$ mit $\overrightarrow{DC} = (3 | 0 | 4)$: parallel; $\overrightarrow{AB} \circ \overrightarrow{AD} = (6 | 0 | 8) \circ (4 | 1 | -3) = 0$: rechter Winkel bei $A$. \quad j) Die Seite an $P$ liegt in der Ebene und steht senkrecht auf $\overrightarrow{PQ} = (0 | 0 | 5)$: Richtung $(4 | -3 | 0)$, ihre Länge $5$ passt schon. $R = P + (4 | -3 | 0) = (6 | -2 | 0)$ (oder $(-2 | 4 | 0)$). \quad k) Die Seite an $P$ liegt in der Ebene und steht senkrecht auf $\overrightarrow{PQ} = (0 | 5 | 0)$: Richtung $(3 | 0 | 4)$, ihre Länge $5$ passt schon. $R = P + (3 | 0 | 4) = (6 | 1 | 8)$ (oder $(0 | 1 | 0)$).}
\erg{15}{a) $B'(3 | b | 0)$ liegt auf der Parallelen zur $x_2$-Achse durch $B$ (nur $x_2$ ändert sich). $(3 | b | 0) = \overrightarrow{AB'}$ und $(3 | b - 7 | 0) = \overrightarrow{CB'}$ sind die beiden Seiten, die in $B'$ zusammenstoßen; Skalarprodukt null heißt, sie stehen senkrecht – rechter Innenwinkel bei $B'$. \quad b) $B'(b | 4 | 0)$ liegt auf der Parallelen zur $x_1$-Achse durch $B$. $(b | 4 | 0) = \overrightarrow{AB'}$ und $(b - 9 | 4 | 0) = \overrightarrow{CB'}$ (Spitze minus Fuß: $B' - C$) sind die Seiten an $B'$; Skalarprodukt null heißt rechter Innenwinkel bei $B'$. \quad c) $P = A + 0{,}36 \cdot \overrightarrow{AB}$ liegt auf der Seite $AB$. Der Inkreismittelpunkt ist der Diagonalenschnittpunkt $M(2 | 3 | 1)$; $\overrightarrow{MP} \circ \overrightarrow{AB} = (1{,}92 | 0 | 1{,}44) \circ (-3 | 0 | 4) = 0$: Der Radius $MP$ steht senkrecht auf der Seite, also berührt der Kreis dort. \quad d) $P = A + 0{,}64 \cdot \overrightarrow{AB}$ liegt auf der Seite $AB$. Der Inkreismittelpunkt ist der Diagonalenschnittpunkt $M(1 | 1 | 2)$; $\overrightarrow{MP} \circ \overrightarrow{AB} = (1{,}92 | 0 | 1{,}44) \circ (3 | 0 | -4) = 0$: Der Radius $MP$ steht senkrecht auf der Seite, also berührt der Kreis dort. \quad e) Die Diagonale $AC$ halbiert die Wand. $P$ liegt zwischen $D$ und $C$, also ist das Dreieck $APD$ kleiner als das Dreieck $ACD$, das die Hälfte ist: Der Teil $ABCP$ ist größer. \quad f) Die Diagonale $BD$ halbiert die Wand. $Q$ liegt zwischen $A$ und $D$, also ist das Dreieck $ABQ$ kleiner als das Dreieck $ABD$, die Hälfte: Der Teil $BCDQ$ ist größer.}
\erg{16}{a) Sie liegt parallel zur $x_1x_2$-Ebene in der Höhe $4$.}
\erg{17}{a) $F$ liegt senkrecht über $C$ in der Höhe von $D$: $F(2 | 4 | 3)$ (nicht mit $E$ verwechseln). \quad b) $\overrightarrow{AD} = \overrightarrow{EH} = (0 | 3 | 0)$, $\overrightarrow{AE} = (0 | 0 | 5)$: $G = B + \overrightarrow{AD} + \overrightarrow{AE} = (6 | 4 | 5)$ (nicht $C$). \quad c) $\overrightarrow{AD} = \overrightarrow{EH} = (0 | 5 | 0)$, $\overrightarrow{AE} = (0 | 0 | 3)$: $G = B + \overrightarrow{AD} + \overrightarrow{AE} = (4 | 7 | 4)$ (nicht $C$). \quad d) $F$ liegt senkrecht über $C$ in der Höhe von $D$: $F(8 | 1 | 7)$ (nicht mit $E$ verwechseln). \quad e) Z. B. rechter Winkel im Ursprung, die Katheten auf die Achsen: $(0 | 0 | 0)$, $(6 | 0 | 0)$, $(0 | 8 | 0)$, Spitze $(0 | 0 | 5)$ – die Hypotenuse $10$ gehört nicht auf eine Achse. \quad f) $C(15 | 15 | 0)$, also $G(15 | 15 | 7)$ und $H(0 | 15 | 7)$. \quad g) Nur die Ecke $(0 | 2 | 2)$ passt: danach $(-1 | 1 | 3)$. \quad h) Nur die Ecke $(4 | 0 | 0)$ passt: danach $(2 | 1 | -1)$. \quad i) Nur die Ecke $(3 | 0 | 2)$ passt: danach $(1 | -3 | 1)$. \quad j) Alle Grundecken haben dieselbe dritte Koordinate: Die Grundfläche ist parallel zur $x_1x_2$-Ebene. Höhe $h = 7 - 2 = 5$ (nicht $7$). \quad k) Alle Grundecken haben dieselbe dritte Koordinate: Die Grundfläche ist parallel zur $x_1x_2$-Ebene. Höhe $h = 3 - (-1) = 4$ (nicht $3$). \quad l) Alle Grundecken haben dieselbe dritte Koordinate: Die Grundfläche ist parallel zur $x_1x_2$-Ebene. Höhe $h = 9 - 3 = 6$ (nicht $9$).}
\erg{18}{a) Strahlensatz am Querschnitt, von der Traufe aus: $d = \frac{2 - 0{,}5}{5{,}5 - 0{,}5} \cdot 4 = 1{,}2$ m (Höhendifferenz, nicht $\frac{2}{5{,}5}$). Nicht zählend: zwei Streifen, Anteil $\frac{2 \cdot 1{,}2}{8} = 30\,\%$. \quad b) Strahlensatz am Querschnitt, von der Traufe aus: $d = \frac{1{,}8 - 1}{5 - 1} \cdot 5 = 1$ m (Höhendifferenz, nicht $\frac{1{,}8}{5}$). Nicht zählend: zwei Streifen, Anteil $\frac{2 \cdot 1}{10} = 20\,\%$. \quad c) Strahlensatz am Querschnitt, von der Traufe aus: $d = \frac{1 - 0{,}4}{2{,}9 - 0{,}4} \cdot 5 = 1{,}2$ m (Höhendifferenz, nicht $\frac{1}{2{,}9}$). Nicht zählend: zwei Streifen, Anteil $\frac{2 \cdot 1{,}2}{10} = 24\,\%$. \quad d) $C$ läuft auf einem Kreis um den Lotfußpunkt $L(4 | 4 | 0)$ auf $AB$ (nicht um den Ursprung), Radius $|\overrightarrow{LC}| = \sqrt{18}$. In der $x_1x_2$-Ebene senkrecht zu $AB$: $L \pm 3 \cdot (1 | 1 | 0)$; mit $x_1 < 4$: $C'(1 | 1 | 0)$. \quad e) $C$ läuft auf einem Kreis um den Lotfußpunkt $L(2 | 5 | 0)$ auf $AB$ (nicht um den Ursprung), Radius $|\overrightarrow{LC}| = 5$. In der $x_1x_2$-Ebene senkrecht zu $AB$: $L \pm 5 \cdot (1 | 0 | 0)$; mit $x_1 > 2$: $C'(7 | 5 | 0)$. \quad f) Ecke $(k | k | k)$ auf $CS$: $C + t \cdot \overrightarrow{CS} = (4 - 4t | 4 - 4t | 6t)$, alle drei Koordinaten gleich: $4 - 4t = 6t$, also $t = 0{,}4$ und $k = 2{,}4$. Die Seitenflächen $SBC$ und $SCD$ liegen in $3x_1 + 2x_3 = 12$ und $3x_2 + 2x_3 = 12$; alle acht Würfelecken erfüllen dort „$\le$“ und liegen damit in der Pyramide – weil beide Körper konvex sind, liegt der ganze Würfel darin (die Prüfung einer Ecke allein genügt nicht). \quad g) Ecke $(k | k | k)$ auf $CS$: $C + t \cdot \overrightarrow{CS} = (3 - 3t | 3 - 3t | 6t)$, alle drei Koordinaten gleich: $3 - 3t = 6t$, also $t = \frac{1}{3}$ und $k = 2$. Die Seitenflächen $SBC$ und $SCD$ liegen in $2x_1 + x_3 = 6$ und $2x_2 + x_3 = 6$; alle acht Würfelecken erfüllen dort „$\le$“ und liegen damit in der Pyramide – weil beide Körper konvex sind, liegt der ganze Würfel darin (die Prüfung einer Ecke allein genügt nicht). \quad h) Für $2 \le s \le 5$ ein Dreieck (drei Ecken): die Bodenstrecke von $x_1 = 2$ bis $x_1 = 6$ und die Spitze auf dem First in der Höhe $3$ – immer dasselbe Dreieck, Flächeninhalt $\frac12 \cdot 4 \cdot 3 = 6$ (Grenzen sind die $x_2$-Koordinaten von $P$ und $Q$). Für $1 < s < 2$ und $5 < s < 7$ ein Trapez (vier Ecken): Die Dreiecksflächen $ABP$ bzw. $CDQ$ schneiden die Spitze ab. \quad i) Für $3 \le s \le 6$ ein Dreieck (drei Ecken): die Bodenstrecke von $x_1 = 0$ bis $x_1 = 6$ und die Spitze auf dem First in der Höhe $4$ – immer dasselbe Dreieck, Flächeninhalt $\frac12 \cdot 6 \cdot 4 = 12$ (Grenzen sind die $x_2$-Koordinaten von $P$ und $Q$). Für $1 < s < 3$ und $6 < s < 9$ ein Trapez (vier Ecken): Die Dreiecksflächen $ABP$ bzw. $CDQ$ schneiden die Spitze ab. \quad j) Die Ebene enthält die Deckendiagonale $EG$ und schneidet den Boden in der dazu parallelen Strecke von $S_t$ bis $(4 | 4 - t | 0)$. Für $0 \le t < 4$ ein Viereck (gleichschenkliges Trapez), für $t = 4$ ein gleichseitiges Dreieck (drei Achsen): $S_{4}$ ist die Ecke $B$, die beiden Bodenpunkte fallen zusammen (nicht dem Viereck zuschlagen). Genau zwei Symmetrieachsen nur für $t = 0$: dann ist es ein Rechteck. \quad k) Die Ebene enthält die Deckendiagonale $EG$ und schneidet den Boden in der dazu parallelen Strecke von $S_t$ bis $(6 | 6 - t | 0)$. Für $0 \le t < 6$ ein Viereck (gleichschenkliges Trapez), für $t = 6$ ein gleichschenkliges Dreieck: $S_{6}$ ist die Ecke $B$, die beiden Bodenpunkte fallen zusammen (nicht dem Viereck zuschlagen). Genau zwei Symmetrieachsen nur für $t = 0$: dann ist es ein Rechteck.}
